\begin{problem}[1 (Schwarzschild with a Cosmological Constant)]
  Consider Einstein's equations in vacuum, but now with a cosmological constant $\Lambda$:
  \[%
    G_{\mu\nu} + \Lambda g_{\mu\nu} = 0
  .\]%
  \begin{enumerate}
    \item (10 points) Solve for the most general, spherically symmetric metric, in coordinates $(t, r, \theta, \phi)$. Show that this reduces to the ordinary Schwarzschild solution when $\Lambda = 0$. You can use information from HW\#5 to minimize the number of computations.
    \item (10 points) Write down the equation of motion for radial geodesics in terms of an effective potential. Sketch the effective potential for massive particles.
  \end{enumerate}
\end{problem}

\begin{solution}[(i)]
  From the results of Homework 5, Problem 3, the most general spherically symmetric metric can be written in the form:
  \[%
    \ds^2 = -e^{\nu(r,t)} \dt^2 + e^{\lambda(r,t)} dr^2 + r^2 \dd{\Omega}^2
  .\]%
  The vacuum Einstein equation $G_{\hat{t}\hat{r}} = 0$ implies that $\lambda$ is independent of time: $\lambda = \lambda(r)$. Furthermore, the combination $G_{\hat{t}\hat{t}} + G_{\hat{r}\hat{r}} = 0$ (which holds here since $G_{\mu\nu} = -\Lambda g_{\mu\nu}$ and $g_{\hat{t}\hat{t}} + g_{\hat{r}\hat{r}} = -1 + 1 = 0$) implies $\nu' + \lambda' = 0$. By choosing a suitable time coordinate, we can set $\nu(r) = -\lambda(r)$.

  Let $e^{-\lambda(r)} = f(r)$. The Einstein component $G_{\hat{t}\hat{t}} = -\Lambda g_{\hat{t}\hat{t}} = \Lambda$. Thus, we get:
  \[%
    G_{\hat{t}\hat{t}} = \frac{e^{-\lambda}(-1 + e^\lambda + r\lambda')}{r^2} = \frac{1}{r^2} \odv{}{r} [r(1 - e^{-\lambda})] = \Lambda
  .\]%
  Integrating with respect to $r$:
  \[%
    r(1 - f(r)) = \int \Lambda r^2 \dr = \frac{\Lambda r^3}{3} + C
  .\]%
  Setting the integration constant $C = 2M$ to match the Newtonian limit (Schwarzschild mass), we find:
  \[%
    f(r) = 1 - \frac{2M}{r} - \frac{\Lambda r^2}{3}
  .\]%
  The most general metric is the Schwarzschild-de Sitter metric:
  \[%
    \ds^2 = -\left(1 - \frac{2M}{r} - \frac{\Lambda r^2}{3}\right) \dt^2 + \left(1 - \frac{2M}{r} - \frac{\Lambda r^2}{3}\right)^{-1} \dr^2 + r^2 \dd{\Omega}^2
  .\]%
  When $\Lambda = 0$, this yields $f(r) = 1 - 2M/r$, which is exactly the ordinary Schwarzschild solution.
\end{solution}

\begin{solution}[(ii)]
  For a massive particle on a radial geodesic ($L=0$), the 4-velocity normalization $g_{\mu\nu}u^\mu u^\nu = -1$ gives:
  \[%
    -f(r) \dot{t}^2 + f(r)^{-1} \dot{r}^2 = -1
  .\]%
  Using the conserved energy per unit mass $E = f(r) \dot{t}$, we substitute $\dot{t} = E/f(r)$:
  \[%
    -\frac{E^2}{f(r)} + \frac{1}{f(r)} \left(\odv{r}{\tau}\right)^2 = -1 \implies \left(\odv{r}{\tau}\right)^2 = E^2 - f(r)
  .\]%
  We define the effective potential $V_{\text{eff}}(r)$ such that $\left(\odv{t}{\tau}\right)^2 + V_{\text{eff}}(r) = E^2$:
  \[%
    V_{\text{eff}}(r) = f(r) = 1 - \frac{2M}{r} - \frac{\Lambda r^2}{3}
  .\]%

  For $\Lambda > 0$, the potential starts at $-\infty$ as $r \to 0$, rises to a local maximum, and then drops back toward $-\infty$ as $r \to \infty$ due to the $-\Lambda r^2$ term.
\end{solution}

\begin{problem}[2 (Metric perturbations)]
  Consider the metric perturbation
  \[%
    h_{\alpha\beta}(x) = \begin{pmatrix}
      c & 0 & 0 & 0 \\
      0 & -c & 0 & 0 \\
      0 & 0 & a & 0 \\
      0 & 0 & 0 & -a \\
    \end{pmatrix} \sin(k)(x - t)
  ,\]%
  where $c, a$ are arbitrary constants.
  \begin{enumerate}
    \item (10 points) Show how this solves the linearized Einstein equation $\delta R_{\alpha\beta} = 0$.
    \item (10 points) Find the functions $\zeta_\alpha(x)$ that transform it into a form where $c = 0$. What is the new value of $a$?
  \end{enumerate}
\end{problem}

\begin{solution}[(i)]
  We consider the linearized Ricci tensor in vacuum ($\delta R_{\alpha\beta} = 0$), which is defined as:
  \[%
    \delta R_{\alpha\beta} = \frac{1}{2} \left(\pd{\alpha} \pd{\mu} h^\mu_\beta + \pd{\beta} \pd{\mu} h^\mu_\alpha - \square h_{\alpha\beta} - \pd{\alpha} \pd{\beta} h\right) = 0
  .\]%
  Let $u = k(x-t)$. The derivatives are $\pd{t} = -k \odv{}{u}$ and $\pd{x} = k \odv{}{u}$. First, we compute the trace $h = \eta^{\mu\nu}h_{\mu\nu} = -h_{00} + h_{11} + h_{22} + h_{33}$:
  \[%
    h = (-c - c + a - a)\sin(u) = -2c\sin(u)
  .\]%
  Since $h_{\alpha\beta}$ depends only on $(x-t)$, the wave equation $\square h_{\alpha\beta} = (-\pd{t}^2 + \pd{x}^2)h_{\alpha\beta} = 0$ is automatically satisfied. Next, we evaluate the divergence terms $\pd{\mu} h^\mu_\alpha$:
  \begin{itemize}
    \item $\pd{\mu} h^\mu_0 = \eta^{00}\pd{t} h_{00} = (-1)(-kc \cos u) = kc \cos u$
    \item $\pd{\mu} h^\mu_1 = \eta^{11}\pd{x} h_{11} = (1)(-kc \cos u) = -kc \cos u$
    \item $\pd{\mu} h^\mu_2 = \pd{\mu} h^\mu_3 = 0$
  \end{itemize}
  Substituting into $\delta R_{00}$:
  \[%
    \delta R_{00} = \frac{1}{2} [ 2\pd{t}(kc \cos u) - \pd{t}^2(-2c \sin u) ] = \frac{1}{2} [ -2k^2c \sin u + 2k^2c \sin u ] = 0
  .\]%
  All other components similarly vanish due to the plane-wave symmetry, confirming the solution.
\end{solution}

\begin{solution}[(ii)]
  A gauge transformation is defined by $h_{\alpha\beta}^{\text{new}} = h_{\alpha\beta} - \pd{\alpha} \zeta_\beta - \pd{\beta} \zeta_\alpha$. To set the new $h_{00} = 0$, we require:
  \[%
    c \sin(k(x-t)) - 2\pd{t} \zeta_0 = 0 \implies \pd{t} \zeta_0 = \frac{c}{2} \sin(k(x-t))
  .\]%
  Integrating with respect to $t$:
  \[%
    \zeta_0 = \frac{c}{2k} \cos(k(x-t))
  .\]%
  Similarly, for $h_{11}^{\text{new}} = 0$:
  \[%
    -c \sin(k(x-t)) - 2\pd{x} \zeta_1 = 0 \implies \pd{x} \zeta_1 = -\frac{c}{2} \sin(k(x-t))
  .\]%
  Integrating with respect to $x$:
  \[%
    \zeta_1 = \frac{c}{2k} \cos(k(x-t))
  .\]%
  With $\zeta_2 = \zeta_3 = 0$, we check the cross term $h_{01}^{\text{new}}$:
  \[%
    h_{01}^{\text{new}} = 0 - \pd{t} \zeta_1 - \pd{x} \zeta_0 = -(\frac{c}{2} \sin u) - (-\frac{c}{2} \sin u) = 0
  .\]%
  Since $\zeta_2 = \zeta_3 = 0$, the transverse components $h_{22}$ and $h_{33}$ remain unchanged. Therefore, the new value of $a$ is:
  \[%
    a_{\text{new}} = a
  .\]%
  This demonstrates that $c$ was a gauge-dependent quantity, while $a$ represents the physical transverse-traceless degrees of freedom.
\end{solution}
